\label{sec:conclusion}

We formulated and tested a conditional equivalence framework for three
placements of automatic differentiation in nonlinear finite elements. Element
and constitutive differentiation agree when they act on the same
quadrature-defined potential, affine constraint map, strain convention, and
smooth active branch. Colored sparse recovery produces the same operator only
under additional pattern, interference, overlap, and row-ownership conditions.

The prerelease diagnostic verification supports element--constitutive
equivalence on smooth local energies and high-order branch-interior plasticity
states. A separate canonical distributed hyperelastic check verifies
rank-local element-Hessian assembly, while distributed colored recovery
remains open. Independent manufactured solutions recover the expected
P1 spatial behavior, including nonaffine hyperelastic displacement,
deformation-gradient, and stress convergence. Conversely, the enriched-rule
study rejects endpoint energy as a sufficient quadrature diagnostic.

These results must be reproduced from the clean release commit before
submission. They diagnose derivative correctness on the tested states, not a
universal computational ranking. The available timings fail the prespecified
publication-model gates, so this revision makes no speedup, crossover, or
scaling claim. The next empirical step is the prepared paired CPU/MPI design
with common Riesz stopping, repeated within-allocation route blocks, canonical
endpoint comparison, and factorized cost measurements. Its outcome may support
a route selector; the mathematical equivalence and evidence hierarchy do not
depend on that outcome.
